% GENERATED FILE. Edit canonical lessons, then run npm run generate:books.
% canonical-source-hash: fnv1a64:215df098246d0346
% canonical-lessons: MR-C127-hivala, MR-C127-pavsala, MR-C127-saral, MR-C127-valan, MR-C127-kopra

\chapter{Winter, Monsoon, Straight, Turn, Corner}
\label{ch:mr-winter-monsoon-straight-turn-corner}
\hlchaptermodality{127}

\begin{chapteropening}
\textbf{By the end of this chapter:} \emph{I can say winter, the monsoon, straight, a turn and a corner.}

Complete the last lesson of chapter 127: I can say winter, the monsoon, straight, a turn and a corner.
\end{chapteropening}

\section[\texorpdfstring{\mr{हिवाळा} (hivāḷā)}{हिवाळा (hivāḷā)}]{\mr{हिवाळा} (hivāḷā) — winter}

\label{lesson:MR-C127-hivala}



\begin{quote}
Before the new one: say the Marathi for every day, then the Marathi for summer.
\end{quote}

\subsection*{You'll want to know: \mr{हिवाळा}}
\textbf{\mr{हिवाळा}} — \emph{hivāḷā} — \textquotedblleft{}winter\textquotedblright{}.

\begin{cousinweb}
\textbf{\mr{हिवाळा}} (\emph{hivāḷā}) is a tadbhava, a word inherited and reshaped by sound change. It comes from \emph{*himakāla}, \textquotedblleft{}winter-time\textquotedblright{}, a form CDIAL reconstructs from its descendants (CDIAL 14098). It is \emph{hima}, \textquotedblleft{}cold, snow\textquotedblright{}, and \emph{kāla}, \textquotedblleft{}time\textquotedblright{}.
\end{cousinweb}

\begin{grammarlens}[title={What you've built}]
One more word for this chapter.
\end{grammarlens}

\subsection*{Your turn}
\begin{itemize}
\raggedright
  \item \emph{Say it:} \emph{hivāḷā}
  \item \emph{Say it:} \emph{hivāḷā}, once more
  \item \emph{Say it:} say \emph{hivāḷā}
  \item \emph{Recall:} say the Marathi for every day, then the Marathi for summer, then say \emph{hivāḷā} again
\end{itemize}

\begin{tcolorbox}[breakable,colback=teal!4,colframe=teal!35!black,title={Before you move on}]
What does \mr{हिवाळा} mean? (\textquotedblleft{}Winter\textquotedblright{}.) Say it once more.
\end{tcolorbox}
\section[\texorpdfstring{\mr{पावसाळा} (pāvsāḷā)}{पावसाळा (pāvsāḷā)}]{\mr{पावसाळा} (pāvsāḷā) — the monsoon}

\label{lesson:MR-C127-pavsala}



\begin{quote}
Before the new one: say the Marathi for summer, then the Marathi for winter.
\end{quote}

\subsection*{You'll want to know: \mr{पावसाळा}}
\textbf{\mr{पावसाळा}} — \emph{pāvsāḷā} — \textquotedblleft{}the monsoon\textquotedblright{}.

\begin{cousinweb}
\textbf{\mr{पावसाळा}} (\emph{pāvsāḷā}) is a tadbhava, a word inherited and reshaped by sound change. It comes from \emph{*prāvṛṣakāla}, \textquotedblleft{}the rainy season\textquotedblright{}, a form CDIAL reconstructs from its descendants (CDIAL 8965). It is \emph{prāvṛṣ}, \textquotedblleft{}the rains\textquotedblright{}, and \emph{kāla}, \textquotedblleft{}time\textquotedblright{}.
\end{cousinweb}

\begin{grammarlens}[title={What you've built}]
One more word for this chapter.
\end{grammarlens}

\subsection*{Your turn}
\begin{itemize}
\raggedright
  \item \emph{Say it:} \emph{pāvsāḷā}
  \item \emph{Say it:} \emph{pāvsāḷā}, once more
  \item \emph{Say it:} say \emph{pāvsāḷā}
  \item \emph{Recall:} say the Marathi for summer, then the Marathi for winter, then say \emph{pāvsāḷā} again
\end{itemize}

\begin{tcolorbox}[breakable,colback=teal!4,colframe=teal!35!black,title={Before you move on}]
What does \mr{पावसाळा} mean? (\textquotedblleft{}The monsoon\textquotedblright{}.) Say it once more.
\end{tcolorbox}
\section[\texorpdfstring{\mr{सरळ} (saraḷ)}{सरळ (saraḷ)}]{\mr{सरळ} (saraḷ) — straight}

\label{lesson:MR-C127-saral}



\begin{quote}
Before the new one: say the Marathi for winter, then the Marathi for the monsoon.
\end{quote}

\subsection*{You'll want to know: \mr{सरळ}}
\textbf{\mr{सरळ}} — \emph{saraḷ} — \textquotedblleft{}straight\textquotedblright{}.

\begin{cousinweb}
\textbf{\mr{सरळ}} (\emph{saraḷ}) is a tadbhava, a word inherited from Sanskrit and reshaped by sound change. It comes from \emph{sarala}, \textquotedblleft{}straight\textquotedblright{} (CDIAL 13252).
\end{cousinweb}

\begin{grammarlens}[title={What you've built}]
One more word for this chapter.
\end{grammarlens}

\subsection*{Your turn}
\begin{itemize}
\raggedright
  \item \emph{Say it:} \emph{saraḷ}
  \item \emph{Say it:} \emph{saraḷ}, once more
  \item \emph{Say it:} say \emph{saraḷ}
  \item \emph{Recall:} say the Marathi for winter, then the Marathi for the monsoon, then say \emph{saraḷ} again
\end{itemize}

\begin{tcolorbox}[breakable,colback=teal!4,colframe=teal!35!black,title={Before you move on}]
What does \mr{सरळ} mean? (\textquotedblleft{}Straight\textquotedblright{}.) Say it once more.
\end{tcolorbox}
\section[\texorpdfstring{\mr{वळण} (vaḷaṇ)}{वळण (vaḷaṇ)}]{\mr{वळण} (vaḷaṇ) — a turn, a bend}

\label{lesson:MR-C127-valan}



\begin{quote}
Before the new one: say the Marathi for the monsoon, then the Marathi for straight.
\end{quote}

\subsection*{You'll want to know: \mr{वळण}}
\textbf{\mr{वळण}} — \emph{vaḷaṇ} — \textquotedblleft{}a turn, a bend\textquotedblright{}.

\begin{cousinweb}
\textbf{\mr{वळण}} (\emph{vaḷaṇ}) is a tadbhava, a word inherited from Sanskrit and reshaped by sound change. It comes from \emph{valana}, \textquotedblleft{}moving in a circle, turning\textquotedblright{} (CDIAL 11406). It is built on \emph{val}, \textquotedblleft{}to turn\textquotedblright{}.
\end{cousinweb}

\begin{grammarlens}[title={What you've built}]
One more word for this chapter.
\end{grammarlens}

\subsection*{Your turn}
\begin{itemize}
\raggedright
  \item \emph{Say it:} \emph{vaḷaṇ}
  \item \emph{Say it:} \emph{vaḷaṇ}, once more
  \item \emph{Say it:} say \emph{vaḷaṇ}
  \item \emph{Recall:} say the Marathi for the monsoon, then the Marathi for straight, then say \emph{vaḷaṇ} again
\end{itemize}

\begin{tcolorbox}[breakable,colback=teal!4,colframe=teal!35!black,title={Before you move on}]
What does \mr{वळण} mean? (\textquotedblleft{}A turn, a bend\textquotedblright{}.) Say it once more.
\end{tcolorbox}
\section[\texorpdfstring{\mr{कोपरा} (koprā)}{कोपरा (koprā)}]{\mr{कोपरा} (koprā) — a corner}

\label{lesson:MR-C127-kopra}



\begin{quote}
Before the new one: say the Marathi for straight, then the Marathi for a turn.
\end{quote}

\subsection*{You'll want to know: \mr{कोपरा}}
\textbf{\mr{कोपरा}} — \emph{koprā} — \textquotedblleft{}a corner\textquotedblright{}.

\begin{cousinweb}
\textbf{\mr{कोपरा}} (\emph{koprā}), \textquotedblleft{}corner\textquotedblright{}, is inherited. CDIAL lists it beside Marathi \emph{kopar}, \textquotedblleft{}elbow\textquotedblright{}, under Sanskrit \emph{kūrpara}, \textquotedblleft{}elbow\textquotedblright{} (CDIAL 3413).
\end{cousinweb}

\begin{grammarlens}[title={What you've built}]
One more word for this chapter.
\end{grammarlens}

\subsection*{Your turn}
\begin{itemize}
\raggedright
  \item \emph{Say it:} \emph{koprā}
  \item \emph{Say it:} \emph{koprā}, once more
  \item \emph{Say it:} say \emph{koprā}
  \item \emph{Recall:} say the Marathi for straight, then the Marathi for a turn, then say \emph{koprā} again
\end{itemize}

\begin{tcolorbox}[breakable,colback=teal!4,colframe=teal!35!black,title={Before you move on}]
What does \mr{कोपरा} mean? (\textquotedblleft{}A corner\textquotedblright{}.) Say it once more.
\end{tcolorbox}
